% GENERATED FILE. Edit canonical lessons, then run npm run generate:books.
% canonical-source-hash: fnv1a64:2229021d0280f8ca
% canonical-lessons: HI-C26-raat, HI-S139-letter-pha

\chapter{Night}
\label{ch:hi-night}
\hlchaptermodality{26}

\begin{chapteropening}
\textbf{By the end of this chapter:} \emph{I can say night in Hindi and explain why \hi{रात} and Latin's nox are not cousins at all.}

Trace rātri $\to$ ratti $\to$ raat, then state the honest twist: unlike \hi{दिन} and diēs, \hi{रात} and nox share no root.
\end{chapteropening}

\section[raat]{\hi{रात} (raat) — \textquotedblleft{}night,\textquotedblright{} and an honest false cousin of Latin's nox}

\label{lesson:HI-C26-raat}



\begin{quote}
Last lesson, \emph{din} turned out to be a real (if distant) cousin of Latin's \emph{diēs}. This lesson is the mirror image: \emph{raat} looks like it should be \emph{nox}'s cousin too — but it genuinely isn't.
\end{quote}

\subsection*{You'll want to know — \hi{रात} — \textquotedblleft{}night\textquotedblright{}}
\textbf{\hi{रात}} (\textbf{raat}) — \textquotedblleft{}\textbf{night}\textquotedblright{} — is the word already hiding inside \textbf{\hi{आधी} \hi{रात}} (\emph{ādhī rāt}, \textquotedblleft{}midnight,\textquotedblright{} Chapter 17). It comes down through a completely ordinary sound-change chain: \textbf{Sanskrit} \textbf{\hi{रात्रि}} (\emph{rātri}) $\to$ \textbf{Prakrit} \textbf{\hi{रत्ति}} (\emph{ratti}, consonant cluster simplifying) $\to$ Hindi \textbf{\hi{रात}} (\emph{raat}) — the same shape of erosion you've already seen in Hindi's numbers (\emph{saptá} $\to$ \emph{sāt}, \emph{aṣṭá} $\to$ \emph{āṭh}).

\subsection*{You'll want to know — A genuine false cousin — not a true cognate}
Here's the honest twist: \textbf{\hi{रात्रि}} (\emph{rātri}) traces to \textbf{Proto-Indo- Iranian} \emph{\textbf{*HráHtriH}}, from \textbf{Proto-Indo-European} \emph{\textbf{*h\textsubscript{1}reh\textsubscript{1}-trih\textsubscript{2}}} — built from a root meaning \textquotedblleft{}\textbf{to rest, remain still},\textquotedblright{} plus a feminine suffix. This is a \textbf{completely different} PIE root from Latin's \textbf{nox} (from \emph{\textbf{*nók\textsuperscript{w}ts}}, the root behind English \textquotedblleft{}night\textquotedblright{} itself, German \emph{Nacht}, Greek \emph{nyx}). Both \emph{rātri} and \emph{nox} are genuine Indo-European \textquotedblleft{}night\textquotedblright{} words, in related language branches — but they are \textbf{not} cousins. Unlike \emph{din} and \emph{diēs} last lesson, which really do share one ancestor, \emph{raat} and \emph{nox} just happen to both mean \textquotedblleft{}night,\textquotedblright{} from two entirely separate ancient roots. Sanskrit actually \textbf{does} keep a true \emph{nók\textsuperscript{w}ts} cousin of its own — \textbf{\hi{नक्तम्}} (\emph{naktam}, \textquotedblleft{}by night\textquotedblright{}) — but it's an archaic, adverbial word; it's \emph{rātri}, not \emph{naktam}, that won the job of everyday Hindi's \emph{raat}.

\subsection*{Your turn}
Say these aloud:

\begin{itemize}
\raggedright
  \item \textquotedblleft{}raat\textquotedblright{} — \textquotedblleft{}night,\textquotedblright{} already met inside ādhī rāt
  \item the sound chain — rātri $\to$ ratti $\to$ raat, the same erosion pattern as saptá $\to$ sāt
  \item the honest twist — raat and nox are NOT cousins, unlike din and diēs
\end{itemize}

\begin{tcolorbox}[breakable,colback=teal!4,colframe=teal!35!black,title={Before you move on}]
Where have you already met \textbf{\hi{रात}} before this lesson? (\textbf{\hi{आधी} \hi{रात}}, \emph{ādhī rāt}, \textquotedblleft{}midnight.\textquotedblright{}) What Sanskrit word, and what sound-change chain, gives Hindi \textbf{raat}? (\textbf{\hi{रात्रि}} \emph{rātri} $\to$ Prakrit \emph{ratti} $\to$ \emph{raat}.) Are \emph{raat} and Latin's \emph{nox} true cognates, the way \emph{din} and \emph{diēs} are? (\textbf{No} — genuinely different PIE roots, \emph{*h\textsubscript{1}reh\textsubscript{1}-} \textquotedblleft{}to rest\textquotedblright{} vs. \emph{*nók\textsuperscript{w}ts} — a false cousin, despite both meaning \textquotedblleft{}night.\textquotedblright{})
\end{tcolorbox}
\section[pha]{\hi{फ} — the fourth breath pair}

\label{lesson:HI-S139-letter-pha}



\begin{quote}
Before the new one: \textbf{\hi{ग}} — what sound does it carry?

Now say \textbf{\hi{प}}, then say it again with a puff of air behind it.
\end{quote}

\subsection*{Script you'll notice: \hi{फ}}
\textbf{\hi{फ}} — \emph{pha}.

It is the \textbf{breathy partner of \hi{प}}, the fourth such pair you have met.

What it is made of:

\begin{itemize}
\raggedright
  \item a left stem joined through the lower bowl to a rising and retraced central stem
  \item a separate clockwise right arch
  \item the top shirorekhā
\end{itemize}

\subsection*{Script: four pairs, one idea}
\noindent
\begin{tabularx}{\linewidth}{@{}>{\raggedright\arraybackslash}X>{\raggedright\arraybackslash}X@{}}
\toprule
\textbf{plain} & \textbf{breathy} \\
\midrule
\hi{त} & \hi{थ} \\
\hi{च} & \hi{छ} \\
\hi{ज} & \hi{झ} \\
\textbf{\hi{प}} & \textbf{\hi{फ}} \\
\bottomrule
\end{tabularx}

Four pairs now. \textbf{The distinction was learned once and each pair since has cost only a shape.}

Like \textbf{\hi{ग}}, this one is late: \textbf{\hi{फिर}} — \emph{again}, in \emph{see you again} — has been yours since the fourth chapter, and the letter opening it went unnamed until now.

\subsection*{Writing: \hi{फ}}
\begin{itemize}
\raggedright
  \item \textbf{1.} descend the left stem, curve right around the lower bowl, rise along the central side to the headline, then retrace down the central stem without lifting
  \item \textbf{2.} lift at the central junction and sweep clockwise over and down through the open right arch
  \item \textbf{3.} lift and draw the top shirorekhā left-to-right
\end{itemize}

\textbf{Pen lifts: 2.} The pen comes up 2 times and no more.

\begin{quote}
Verified three-stroke teaching form; another learner source stages the components differently.
\end{quote}

\begin{quote}
This is one attested teaching order and not a national standard — handwriting here is taught with school-to-school variation. Source: JackPotte, ‘Devanagari p\textsuperscript{h} \hi{फ}.gif’, strokes 1–3, Wikimedia Commons, 29 March 2009.
\end{quote}

\subsection*{Your turn}
\begin{itemize}
\raggedright
  \item \emph{Look:} at \hi{प} and \hi{फ} side by side, and say which one carries the puff
\end{itemize}

\begin{quote}
\hi{प}  ·  \hi{फ}
\end{quote}

\begin{itemize}
\raggedright
  \item \emph{Trace it:} \hi{फ} three times, saying \emph{pha} as you finish each one
  \item \emph{Say it:} the three earlier pairs that work the same way
\end{itemize}

\begin{tcolorbox}[breakable,colback=teal!4,colframe=teal!35!black,title={Before you move on}]
Which letter is \textbf{\hi{फ}} the breathy partner of? (\textbf{\hi{प}}.) How many breath pairs have you met? (Four.) Name a word you have said since early on that begins with it. (\textbf{\hi{फिर}}.)
\end{tcolorbox}
